% p1_evaluation_protocol.tex
% Phase 1 — Evaluation Protocol and Theory
% Author: Krishna Pramod Dubale (Task 04)
% ──────────────────────────────────────────────────────────────────────────────

\subsection{Evaluation Metrics Under Class Imbalance}
\label{sec:eval:metrics}

With a fraud base rate of $0.17\%$, trivial majority prediction achieves $99.83\%$ accuracy while failing all fraud detections. We strictly report imbalanced metrics:
\begin{itemize}
    \item \textbf{PR-AUC (Average Precision):} The primary headline metric, computing the area under the precision-recall curve $\mathrm{AP} = \sum_k (R_k - R_{k-1}) P_k$. Unlike ROC-AUC, PR-AUC is invariant to true negative inflation and directly penalizes false alarms.
    \item \textbf{Recall at Fixed FPR:} In banking operations with constrained false-alarm budgets, we report $\text{Recall@FPR} \le x$ ($x \in \{0.1\%, 0.5\%, 1.0\%\}$). Thresholds are derived from the $(1-x)$-quantile of legitimate validation scores, frozen, and evaluated forward on test.
    \item \textbf{Precision at Capacity:} Measuring investigator efficiency under fixed queue limits: $\text{Precision@}K = \text{Frauds in Top-}K / K$ for $K \in \{50, 100, 200\}$.
\end{itemize}

\subsection{Stratified Bootstrap Confidence Intervals}
To quantify sampling uncertainty over the small number of test fraud events (\NFraudTest{}), we compute 95\% confidence intervals via stratified bootstrap with $B = 2\,000$ resamples. Resampling legitimate and fraud instances separately preserves empirical prevalence across replicates. CIs are determined via the 2.5th and 97.5th percentiles of the bootstrap metric distribution. Thresholds calibrated on validation data are applied out-of-sample to ensure zero test-set leakage.
